% =============================================================================
% Presentación Beamer P03 – Anemometría de Hilo Caliente
% Técnicas de Medida – Ingeniería Aeronáutica
% =============================================================================
\documentclass[aspectratio=169,12pt]{beamer}

\usetheme{Madrid}
\usecolortheme{seahorse}
\usefonttheme{professionalfonts}

\usepackage[utf8]{inputenc}
\usepackage[T1]{fontenc}
\usepackage[spanish]{babel}
\usepackage{amsmath,amssymb}
\usepackage{siunitx}
\sisetup{output-decimal-marker={,}}
\usepackage{tikz}
\usetikzlibrary{arrows.meta,positioning,calc,shapes.geometric}
\usepackage{circuitikz}
\usepackage{booktabs}
\usepackage{listings}
\lstset{
  basicstyle=\ttfamily\tiny,
  keywordstyle=\color{blue!80}\bfseries,
  commentstyle=\color{green!50!black},
  backgroundcolor=\color{gray!10},
  frame=single, breaklines=true
}

\definecolor{aeroblue}{RGB}{0,51,102}
\definecolor{aerored}{RGB}{153,0,0}
\definecolor{aerogreen}{RGB}{0,102,51}
\definecolor{aeroorange}{RGB}{200,100,0}

\setbeamercolor{frametitle}{bg=aeroblue,fg=white}
\setbeamercolor{title}{bg=aeroblue,fg=white}
\setbeamercolor{structure}{fg=aeroblue}

\title[P03 -- Hilo Caliente]{P03: Anemometría de Hilo Caliente (CTA/HWA)}
\subtitle{Técnicas de Medida -- Ingeniería Aeronáutica}
\author{Laboratorio de Técnicas de Medida}
\institute{Departamento de Ingeniería Aeroespacial}
\date{\today}

% =============================================================================
\begin{document}

\begin{frame}
  \titlepage
\end{frame}

\begin{frame}{Contenido}
  \tableofcontents
\end{frame}

% =============================================================================
\section{Principio de Funcionamiento}

\begin{frame}{¿Qué es el HWA?}
  \begin{columns}
    \begin{column}{0.55\textwidth}
      \textbf{Hot-Wire Anemometry (HWA / CTA):}
      \begin{itemize}
        \item Hilo de \textbf{wolframio}, $d \approx \SI{5}{\micro\m}$,
              $L \approx \SI{1}{\mm}$
        \item Calentado eléctricamente ($T_w \approx \SI{220}{\celsius}$)
        \item El flujo \textbf{enfría el hilo} $\to$ varía la resistencia
        \item Puente CTA mantiene $T_w = \mathrm{cte}$ ajustando la corriente
        \item $E(t) \leftrightarrow U(t)$ vía \textbf{ley de King}
      \end{itemize}
      \vspace{0.3em}
      \begin{alertblock}{Ventaja clave}
        Ancho de banda hasta \SI{400}{\kHz} -- ideal para turbulencia.
      \end{alertblock}
    \end{column}
    \begin{column}{0.42\textwidth}
      \begin{tikzpicture}[scale=0.85,>=Stealth]
        % Probe body
        \draw[fill=gray!30,draw=black,thick] (1.5,-0.4) rectangle (1.9,1.8);
        \draw[fill=gray!50] (1.5,1.8) -- (1.7,2.0) -- (1.9,1.8);
        % Prongs
        \draw[black,very thick] (1.7,2.0) -- (1.2,2.5);
        \draw[black,very thick] (1.7,2.0) -- (2.2,2.5);
        % Hot wire
        \draw[aerored,ultra thick,line width=2.5pt] (1.2,2.5) -- (2.2,2.5);
        \node[aerored,above,font=\tiny\bfseries] at (1.7,2.6) {Hilo ($d=5\,\mu$m)};
        % Flow
        \foreach \y in {2.2,2.5,2.8} {
          \draw[->,blue!60,thick] (0.2,\y) -- (1.1,\y);
        }
        \node[blue!60,font=\tiny] at (0.6,2.0) {$U$};
        % Heat arrows
        \draw[aerored!60,->] (1.7,2.5) -- (1.7,3.2)
              node[above,font=\tiny,aerored] {$Q_\mathrm{conv}$};
        % Connection
        \draw[aeroblue,thick,->] (1.9,1.0) -- (3.0,1.0)
              node[right,font=\tiny] {$E(t)$};
        % CTA label
        \node[aeroblue,font=\tiny,right] at (3.0,1.3) {CTA};
      \end{tikzpicture}
    \end{column}
  \end{columns}
\end{frame}

% =============================================================================
\section{Teoría: Ley de King}

\begin{frame}{Ley de King}
  \begin{columns}
    \begin{column}{0.5\textwidth}
      \textbf{Balance energético:}
      \begin{equation*}
        I_w^2 R_w = h\,A_s\,(T_w - T_a)
      \end{equation*}

      \textbf{Nusselt cilindrico:} $Nu = A_0 + B_0 Re^n$

      \vspace{0.5em}
      \begin{block}{Ley de King}
        \begin{equation*}
          \boxed{E^2 = A + B\,U^n}
        \end{equation*}
        $A$, $B$ constantes del sensor\\
        $n \approx 0.45$ (King) o ajustado
      \end{block}

      \vspace{0.3em}
      \textbf{Inversión:}
      \begin{equation*}
        U = \left(\frac{E^2 - A}{B}\right)^{1/n}
      \end{equation*}
    \end{column}
    \begin{column}{0.47\textwidth}
      \begin{tikzpicture}[scale=0.75]
        \draw[->] (0,0) -- (5,0) node[right,font=\tiny] {$U$ (m/s)};
        \draw[->] (0,0) -- (0,4) node[above,font=\tiny] {$E^2$ (V²)};
        % King curve (E^2 = 0.9 + 2.5*U^(1/3))
        \draw[aerored,very thick] plot[domain=0:4.5,samples=50]
              (\x, {0.23 + 0.63*\x^0.333});
        % Data points
        \foreach \u/\e in {0.5/0.46, 1.0/0.63, 1.5/0.78, 2.0/0.84,
                            2.5/0.92, 3.5/1.06, 4.5/1.28} {
          \fill[aeroblue] (\u,\e) circle (0.07);
        }
        \node[aerored,font=\tiny] at (3.5,0.8) {$E^2=A+BU^n$};
        % Axes tick labels
        \foreach \x/\l in {1/10,2/20,3/30,4/40} {
          \node[below,font=\tiny] at (\x,0) {\l};
        }
        \foreach \y/\l in {0.5/4.6,1.0/6.3,1.5/8.4,2.0/10.6} {
          \node[left,font=\tiny] at (0,\y) {\l};
        }
      \end{tikzpicture}
    \end{column}
  \end{columns}
\end{frame}

\begin{frame}{Circuito CTA -- Puente Wheatstone}
  \begin{columns}
    \begin{column}{0.55\textwidth}
      El puente Wheatstone mantiene $R_w = \mathrm{cte}$:
      \vspace{0.3em}
      \begin{itemize}\small
        \item Si $U\uparrow$ $\Rightarrow$ enfría el hilo
        \item $R_w\downarrow$ $\Rightarrow$ $\Delta V \neq 0$
        \item Amplificador corrige la corriente $I_w\uparrow$
        \item Nuevo equilibrio: $T_w = \mathrm{cte}$
        \item $E_\mathrm{out} \propto I_w \propto U$
      \end{itemize}
      \vspace{0.3em}
      \begin{block}{Resultado}
        $E(t)$ contiene la información de $U(t)$ a alta frecuencia.
      \end{block}
    \end{column}
    \begin{column}{0.42\textwidth}
      \begin{circuitikz}[scale=0.55, transform shape]
        \node (A) at (0,3) {};
        \node (B) at (0,-3) {};
        \node (L) at (-3,0) {};
        \node (R) at (3,0) {};
        \draw (A) to[R,l=$R_1$] (L);
        \draw (A) to[R,l=$R_2$] (R);
        \draw (L) to[R,l=$R_3$] (B);
        \draw[aerored,very thick] (R) to[R,l_=$R_w$,color=aerored] (B);
        \draw (A) to[battery1,l=$V_s$] (0,4.5);
        \draw (0,-3) to[short] (0,-4.5);
        \draw[fill=aeroblue!10] (-2,0.8) -- (2,0.8) -- (0,2.2) -- cycle;
        \node[font=\tiny,aeroblue] at (0,1.4) {Amp.};
        \draw[aeroblue,->] (0,2.2) -- (0,3.3)
              node[above,font=\tiny] {$E_{out}$};
        \draw[aeroorange,dashed,->] (0,3.0) -- (4.5,3.0) -- (4.5,-1) -- (3.2,-0.2);
        \node[aeroorange,font=\tiny] at (4.8,1) {FB};
      \end{circuitikz}
    \end{column}
  \end{columns}
\end{frame}

% =============================================================================
\section{Calibración}

\begin{frame}{Calibración -- Datos del laboratorio}
  \begin{columns}
    \begin{column}{0.5\textwidth}
      \textbf{Datos de calibración (lab.):}
      \begin{table}
        \centering
        \tiny
        \begin{tabular}{SSS}
          \toprule
          {$U$ (m/s)} & {$E$ (V)} & {$E^2$ (V²)} \\
          \midrule
           5.2 & 2.15 & 4.622 \\
          10.4 & 2.52 & 6.350 \\
          15.6 & 2.74 & 7.508 \\
          20.8 & 2.90 & 8.410 \\
          26.0 & 3.03 & 9.181 \\
          36.4 & 3.26 & 10.628 \\
          \bottomrule
        \end{tabular}
      \end{table}
      \textbf{Resultado del ajuste ($n=3$):}
      \begin{align*}
        A &= 0.9141\,\si{V^2}\\
        B &= 2.5486\,\si{V^2}\\
        R^2 &= 0.872
      \end{align*}
    \end{column}
    \begin{column}{0.47\textwidth}
      \begin{tikzpicture}[scale=0.8]
        \draw[->] (0,0) -- (4.5,0) node[right,font=\tiny] {$U^{1/3}$};
        \draw[->] (0,0) -- (0,4.5) node[above,font=\tiny] {$E^2$ (V²)};
        % Linear fit line
        \draw[aeroblue,thick] (1.5,0.5) -- (3.5,3.8);
        % Data points
        \foreach \x/\y in {1.732/1.04, 2.183/1.43, 2.503/1.69,
                            2.748/1.89, 2.962/2.07, 3.314/2.39} {
          \fill[aerored] (\x-0.7,{\y*1.5}) circle (0.1);
        }
        \node[aeroblue,font=\tiny,rotate=52] at (2.8,2.5)
              {$E^2=0.914+2.549\,U^{1/3}$};
        \node[left,font=\tiny] at (0,1.04) {4.6};
        \node[left,font=\tiny] at (0,2.39) {10.6};
        \node[below,font=\tiny] at (1.03,0) {1.73};
        \node[below,font=\tiny] at (2.6,0) {2.96};
      \end{tikzpicture}
    \end{column}
  \end{columns}
\end{frame}

% =============================================================================
\section{Estadísticos de Turbulencia}

\begin{frame}{Turbulencia: estadísticos y espectro}
  \begin{columns}
    \begin{column}{0.5\textwidth}
      \textbf{Velocidad media y fluctuación:}
      \begin{align*}
        \bar{U} &= \SI{26.0}{\m\per\s}\\
        u'_\mathrm{rms} &= \SI{1.82}{\m\per\s}\\
        I_u &= \frac{1.82}{26.0}\times100\% = 7.0\%
      \end{align*}
      \vspace{0.3em}
      \textbf{Ley de Kolmogorov:}
      \begin{equation*}
        S(f) \propto f^{-5/3}
      \end{equation*}
      (sub-rango inercial)
    \end{column}
    \begin{column}{0.47\textwidth}
      \begin{tikzpicture}[scale=0.8]
        \draw[->] (0,0) -- (5,0) node[right,font=\tiny] {$\log f$};
        \draw[->] (0,0) -- (0,4) node[above,font=\tiny] {$\log S(f)$};
        \fill[aeroblue!15] (0,0) rectangle (1.5,4);
        \fill[aerogreen!15] (1.5,0) rectangle (3.5,4);
        \fill[aerored!10] (3.5,0) rectangle (5,4);
        \node[font=\tiny,aeroblue,rotate=90] at (0.75,2) {Energía};
        \node[font=\tiny,aerogreen,rotate=90] at (2.5,2) {Inercial};
        \node[font=\tiny,aerored,rotate=90] at (4.25,2) {Disipac.};
        \draw[aerored,very thick] (1.5,3.5) -- (3.5,0.8);
        \node[aerored,font=\tiny,rotate=-67] at (2.8,2.3) {$-5/3$};
        \draw[aeroblue,thick] (0.5,3.8) .. controls (1.0,3.6) .. (1.5,3.5);
        \draw[aeroblue,thick] (3.5,0.8) -- (4.8,0.2);
      \end{tikzpicture}
    \end{column}
  \end{columns}
\end{frame}

% =============================================================================
\section{Caso de Estudio}

\begin{frame}{Caso: Capa límite sobre NACA 4412}
  \begin{block}{Escenario}
    Perfil de velocidades HWA en capa límite de NACA~4412 ($c=\SI{200}{\mm}$,
    $\alpha=\SI{4}{\degree}$, $Re_c=3\times10^5$).
  \end{block}
  \begin{columns}
    \begin{column}{0.5\textwidth}
      \textbf{Ley logarítmica (turbulenta):}
      \begin{equation*}
        U^+ = \frac{1}{\kappa}\ln(y^+) + B,
        \quad \kappa = 0.41
      \end{equation*}
      \textbf{Grosor de desplazamiento:}
      \begin{equation*}
        \delta^* = \int_0^\infty \left(1-\frac{U}{U_e}\right)dy
      \end{equation*}
      \begin{itemize}\small
        \item Resolución del traverso: \SI{0.5}{\mm}
        \item Barrido: $y \in [0, 20]\,\si{\mm}$
        \item $\delta^* \approx \SI{3.5}{\mm}$
      \end{itemize}
    \end{column}
    \begin{column}{0.47\textwidth}
      \begin{tikzpicture}[scale=1.0,>=Stealth]
        \draw[->] (0,0) -- (0,4) node[above,font=\tiny] {$y$ (mm)};
        \draw[->] (0,0) -- (4,0) node[right,font=\tiny] {$U$ (m/s)};
        % Wall
        \fill[gray!50] (-0.2,-0.2) rectangle (0,4.2);
        % Log law profile
        \draw[aerored,very thick] plot[domain=0.1:3.8,samples=30]
              ({0.8*ln(\x/0.1+1)}, \x);
        % Reference velocity
        \draw[gray,dashed] (2.8,0) -- (2.8,4);
        \node[gray,right,font=\tiny] at (2.8,3.8) {$U_e$};
        % Delta star
        \fill[aeroblue!20] (0,0) -- (2.8,0) -- (2.8,0.7) ..
              controls (1.5,1.5) .. (0,3.5) -- (0,0);
        \node[aeroblue,font=\tiny] at (0.8,1.0) {$\delta^*$};
        \foreach \y/\u in {0.5/0.9, 1.0/1.5, 1.5/2.0, 2.0/2.3,
                            2.5/2.5, 3.0/2.7, 3.5/2.8} {
          \fill[aerored] (\u,\y) circle (0.06);
        }
      \end{tikzpicture}
    \end{column}
  \end{columns}
\end{frame}

% =============================================================================
\section{Notebook Jupyter}

\begin{frame}[fragile]{P03\_Hilo\_Caliente.ipynb}
  \begin{columns}
    \begin{column}{0.55\textwidth}
      \textbf{Ruta:}\\
      \texttt{\small Practicas/P03\_Hilo\_Caliente/\\notebooks/P03\_Hilo\_Caliente.ipynb}

      \vspace{0.5em}
      \textbf{El notebook/script automatiza:}
      \begin{enumerate}\small
        \item Lectura de \texttt{files/p3/} (señales E(t))
        \item Descarte de transitorios (10\% c/extremo)
        \item Calibración King ($A$, $B$, $n$)
        \item $\bar{U}$, $u'_\mathrm{rms}$, $I_u$, PSD
        \item Gráficas y exportación CSV/JSON
        \item Actualización de \texttt{p3.md}
      \end{enumerate}
    \end{column}
    \begin{column}{0.42\textwidth}
      \begin{lstlisting}[language=Python]
from scipy.optimize import curve_fit
import numpy as np

def king(U, A, B, n):
    return A + B * U**n

U_ref = np.array([5.2, 10.4,
                  15.6, 20.8,
                  26.0, 36.4])
E_sq  = np.array([4.622, 6.350,
                  7.508, 8.410,
                  9.181, 10.628])

popt, _ = curve_fit(
    king, U_ref, E_sq,
    p0=[1, 2.5, 0.45])
A, B, n = popt
# A=0.914, B=2.549, n=0.333
      \end{lstlisting}
    \end{column}
  \end{columns}
\end{frame}

% =============================================================================
\section{Resumen}

\begin{frame}{Conclusiones}
  \begin{enumerate}
    \item El HWA con CTA tiene ancho de banda hasta \SI{400}{\kHz} --
          idóneo para turbulencia aeronáutica.
    \item La \textbf{Ley de King} ($E^2 = A+BU^n$) calibra el sensor:
          $A=0.914$, $B=2.549$, $n=1/3$.
    \item $I_u \approx 7\%$ en la sección de prueba a \SI{26}{\m\per\s}.
    \item El PSD muestra la ley $f^{-5/3}$ de Kolmogorov en el sub-rango inercial.
    \item El notebook y el script Python automatizan todo el flujo de trabajo.
  \end{enumerate}
  \vspace{0.5em}
  \begin{block}{Recursos}
    Manual: \texttt{manuals/P03\_Hilo\_Caliente/manual\_P03.tex}\\
    Notebook: \texttt{Practicas/P03\_Hilo\_Caliente/notebooks/P03\_Hilo\_Caliente.ipynb}\\
    Script: \texttt{Practicas/P03\_Hilo\_Caliente/src/P03\_Hilo\_Caliente.py}
  \end{block}
\end{frame}

\begin{frame}{Referencias}
  \begin{itemize}
    \item Bruun, H.H. (1995). \textit{Hot-Wire Anemometry}, Oxford University Press
    \item J{\o}rgensen, F.E. (2002). \textit{How to Measure Turbulence with HWA}, Dantec
    \item Pope, S.B. (2000). \textit{Turbulent Flows}, Cambridge University Press
    \item Kolmogorov, A.N. (1941). Local structure of turbulence, \textit{Dokl. SSSR}
  \end{itemize}
\end{frame}

\end{document}
